% Folder 147, batch 02 (archivists' pages 21-40).
% Pass: Opus 5.5 (claude-opus-5-5), 2026-10-09 — first pass, unchecked against the pages by a human.
%
% Notation fixed for this batch: \mathfrak{X} for his gothic X (the relative
% curve), \mathcal{S} for the base, S (roman) for the vertex set of an
% MD-graph, \underline{I} where he underlines I, \mathfrak{S} for symmetric
% groups.
\documentclass[11pt,a4paper]{article}
\input{../preamble/grothendieck}

\folder{147}
\batch{2}
\pages{21}{40}
\foldertitle{Structure à l'infini des Mg,ν [pages 1 à 90, dont table des matières] : notes manuscrites (s.d.).}
\dating{s.d. — le groupe « Autour de La "Longue Marche" à travers la théorie de Galois » (141 à 149) est daté [à partir de 1978]-1983}
\watermark{Édition de démonstration}

\begin{document}

\page{21}\note{p.~10 de l'auteur. La page s'ouvre sur un NB qui prolonge l'argument du lot précédent, sur les $M_{g,J}$, $M_G$ et $M_{[G]}$.}
\marginal{$(g, \nu = \operatorname{card} J)$) \ill{}}
\emph{NB} Si $\operatorname{card} J = \nu$, alors $M_{g,J}$ est isomorphe au schéma modulaire $\overset{!}{M}_{g,\nu}$ (via un iso.\ $[1,\nu] \cap \mathbb{Z} \simeq J$), de façon compatible avec les automor.\ des groupes $\mathfrak{S}_J$ et $\mathfrak{S}_\nu$.\note{le « ! » est écrit au-dessus du $M$, ici et à la ligne suivante.} On sait que $\overset{!}{M}_{g,\nu}$ est lisse de type fini sur $\mathbb{Z}$. Il en est donc de même des $M_{g,J}$, donc aussi de $M_G$ (pour $G$ semi-stable) et des $M_{[G]} = (M_G, \Gamma)$.

\section{4. La théorie de Mumford-Deligne}

Soient $\mathcal{S}$ un\supplied{e} multiplicité schématique, $\mathfrak{X}$ un \struck{courbe} schéma relatif sur $\mathcal{S}$, propre sur $\mathcal{S}$, $\underline{I}$ un sous-schéma fermé de $\mathfrak{X}$. On dit que $(\mathfrak{X}, \underline{I})$ est une \emph{MD-courbe relative} sur $\mathcal{S}$, si $\mathfrak{X}$, $\underline{I}$ sont plats de présentation finie sur $\mathcal{S}$, et si pour tt pt géom.\ de $\mathcal{S}$, la

\page{22}
fibre $(\mathfrak{X}_{\bar{s}}, \underline{I}_{\bar{s}})$ est une \add{MD-}courbe géométrique sur $k(s)$ --- i.e.\ $\mathfrak{X}_{\bar{s}}$ est connexe, de dim~1, \struck{mais} que ses pts singuliers \add{$\notin \underline{I}_{\bar{s}}$ et sont} des pts doubles ordinaires, et enfin la condition de rigidité infinitésimale (\ill{}) est satisfaite --- \struck{donc} i.e.\ on exige que $\mathfrak{X}_{\bar{s}}$ ne soit pas de genre~1 lisse avec $\underline{I}_{\bar{s}} = \emptyset$, \struck{et} que les comp.\ irréd.\ de genre~0 de $\mathfrak{X}_{\bar{s}}$ aient des « points marqués » $\geqslant 3$.

La différence avec les considérations du n\textsuperscript{o} précédent, c'est que maintenant (sur une base $\mathcal{S}$ connexe, p.ex.\ \struck{\ill{}} un trait) les \emph{types} des fibres géométriques peuvent varier d'un pt géom.\ à un autre. Mais il est clair que le type \emph{numérique} $(g,\nu)$ ne change pas (invariance des $\chi(\hat{X}_s, \underline{O}_{\hat{X}_s})$… et de $\operatorname{card} \underline{I}_s$), plus précisément,

\page{23}
c'est une fonction localement constante sur $\mathcal{S}$.

Fixons nous un type numérique \add{\emph{anabélien} ($2g+\nu \geqslant 3$)} $(g,\nu)$, et considérons, pour $\mathcal{S}$ variable, le groupoïde fibré
\[
\text{(24)}\qquad \mathcal{S} \longmapsto \hat{M}_{g,\nu}(\mathcal{S}) = \text{MD-courbes relatives sur } \mathcal{S}\text{, de type numérique égal à } (g,\nu)
\]
\note{le numéro « (24) » est lu sur un « 4 » ouvert ; la page suivante porte (25).}
On a alors le vraiment magnifique théorème suivant :

\marginal{On suppose $2g+\nu \geqslant 3$ (cas anabélien)}
\emph{Théorème de Mumford-Deligne} : Le groupoïde fibré $\mathcal{S} \mapsto \hat{M}_{g,\nu}(\mathcal{S})$ sur $(\mathrm{Sch})$ (plus gén\supplied{éralemen}t, sur les multiplicités schématiques…) est représentable par une multiplicité schématique $\hat{M}_{g,\nu}$, qui est \emph{lisse} et \emph{propre} sur $\operatorname{Spec} \mathbb{Z}$, \struck{à fibres géom.\ connexes.} \add{D'autre part $M_{g,\nu}$ est un}\note{l'énoncé est encadré d'un trait vertical dans la marge gauche, qui se poursuit en haut de la page suivante.}

\page{24}
ouvert de Zariski de $\hat{M}_{g,\nu}$, schématiquement dense fibre par fibre.

On en déduit aisément p.ex.\ la connexité des fibres géom.\ de $M_{g,\nu}$ en toute car., à partir de \struck{la} connexité en car.~0. D'autre part, pour tout type $[G]$ de MD-graphes, ayant le type numérique $(g,\nu)$, les courbes standard de type $[G]$ (sur une base quelc.\ $\mathcal{S}$) sont \uncertain{en particulier} des MD-courbes relatives de type numérique $(g,\nu)$. On en conclut un morphisme canonique
\[
\text{(25)}\qquad M_{[G]} \hookrightarrow \hat{M}_{g,\nu}
\]
dont on vérifie aisément que c'est une \emph{immersion}. Le cas où $G$ est le

\page{25}\note{p.~12 de l'auteur.}
MD-graphe trivial, ayant un seul sommet de genre $g$, de poids marqué $\nu$, et pas d'arêtes, redonne l'immersion ouverte canonique
\[
\text{(26)}\qquad M_{g,\nu} \hookrightarrow \hat{M}_{g,\nu} .
\]
Identifiant les $M_{[G]}$ à des sous-multiplicités de $\hat{M}_{g,\nu}$, on voit que les $M_{[G]}$ (pour tous les types de \struck{MD}-graphes possibles, de type numérique $(g,\nu)$) forment une \emph{partition} de $\hat{M}_{g,\nu}$ en un ens fini de sous-multiplicités.

\emph{Remarque.} Considérons l'adhérence d'un $M_{[G]}$ dans $\hat{M}_{g,\nu}$, on trouve que c'est une réunion finie \struck{de} des $M_{[G']}$. On dit que $[G']$ est spécialisation de $[G]$ si $M_{G'} \subset \overline{M_G}$ dans $\hat{M}_{g,\nu}$.

\page{26}
C'est là une \emph{relation d'ordre} dans l'ens des classes d'iso.\ de MD-graphes de type numérique $(g,\nu)$, ayant un plus grand élément. Dans la suite, nous allons préciser considérablement cette notion de spécialisation entre MD-graphes, en la reliant à une notion de \emph{morphismes} (de « générisation ») de MD-graphes. …

Notons que chaque $M_{[G]}$ \struck{est isomorphe} \add{\uncertain{est quotient}} d'un revêtement fini $M_G$, qui est représentable comme produit fini
\[
M_G \simeq \prod_\alpha \underbrace{M_{g_\alpha, \hat{I}_\alpha}}_{\dim\, 3(g_\alpha - 1) + \hat{\nu}_\alpha}
\]
donc
\begin{align*}
\text{(27)}\quad \dim M_G/\mathbb{Z}
  &= \sum_{\alpha\in S} 3(g_\alpha - 1) + \sum \hat{\nu}_\alpha \\
  &= 3\sum g_\alpha - 3\underbrace{\operatorname{card} S}_{s_0} + \sum_{\alpha\in S_1} \nu_\alpha + \sum_{\alpha\in S_1} \tilde{\mu}_\alpha \\
  &= 3\sum g_\alpha - 3s_0 + \nu + 2s_1
\end{align*}
\note{avant « (27) », un mot biffé illisible ; sous $\nu$ et $2s_1$ de la dernière ligne, il écrit en petit « poids marqué » et un mot que l'on peut lire \uncertain{ordre}. L'indice de la seconde somme est en partie caché ; on le lit $\alpha \in S_1$ comme celui de la première.}

\page{27}\note{p.~13 de l'auteur.}
où $s_0$, $s_1$ sont resp.\ les nbs de sommets et d'arêtes du graphe $G$. Or
\[
-3s_0 + 2s_1 = -3\underbrace{(s_0 - s_1)}_{1-h_1} - s_1 = -3 + 3h_1 - s_1 = 3(h_1 - 1) - s_1,
\]
d'où
\[
\dim M_G/\mathbb{Z} = 3\bigl(\underbrace{(\textstyle\sum g_\alpha) + h_1}_{g} - 1\bigr) + \nu - s_1
\]
et enfin la formule
\[
\text{(28)}\qquad \dim M_G/\mathbb{Z} = \underbrace{\bigl[3(g-1)+\nu\bigr]}_{= \dim \hat{M}_{g,\nu}/\mathbb{Z}} - s_1 ,
\]
i.e.
\[
\text{(29)}\qquad \operatorname{codim}_{\hat{M}_{g,\nu}} M_{[G]} = s_1
\]
(nb d'arêtes du MD-graphe $G$, i.e.\ « nb de pts doubles » des courbes standard du type $G$).

L'entier (27) est appelé \emph{dimension modulaire} du MD-graphe $G$
\begin{align*}
\text{(30)}\quad \mu(G) &= \sum_{\alpha\in S} \bigl(3(g_\alpha - 1) + \hat{\nu}_\alpha\bigr) = 3\sum g_\alpha - 3s_0 + \nu + 2s_1 \\
  &= \bigl(3(g-1) + \nu\bigr) - s_1
\end{align*}
l'entier $s_1 = \operatorname{card} A = (\operatorname{card} \tilde{A})/2$ s'appelle la codimension modulaire. Elle est nulle ssi $G$ est un MD-graphe trivial réduit : un pt isolé, sans arêtes, mais

\page{28}
avec un poids marqué \struck{\ill{}}, (et bien sûr, un genre donné). \struck{Du} Du pt de vue structure de graphe muni d'un type numérique $(g,\nu)$, c'est le plus simple de tous ! On étudiera les plus complexes tantôt, qui correspondent au cas d'une dimension modulaire nulle, \struck{ou} et d'une codimension modulaire maximale, égale à $3(g-1)+\nu$.

Il est tentant d'introduire le compactifié
\[
\hat{M}_G \overset{\text{déf}}{=} \prod_{\alpha\in S} \hat{M}_{g_\alpha, \hat{I}_\alpha}
\]
où les $\hat{M}_{g_\alpha, \hat{I}_\alpha}$ sont les compactifiés de MD des $M_{g_\alpha, \hat{I}_\alpha} \simeq M_{g_\alpha, \hat{\nu}_\alpha}$. \struck{Ils} \add{Pour $\alpha$ fixé, $\hat{M}_{g_\alpha, \hat{I}_\alpha}$} représente\struck{nt} le groupoïde fibré
\[
\text{(30)}\qquad \mathcal{S} \longmapsto
\begin{array}{l}
\text{courbes de MD relatives } Y_\alpha/\mathcal{S}\text{,} \\
\text{de type } g_\alpha, \hat{I}_\alpha \text{ i.e.\ de genre } g_\alpha\text{, munies d'un mono} \\
(\hat{I}_\alpha)_{\mathcal{S}} \hookrightarrow Y_\alpha^{\text{lisse}}
\end{array}
\]
\note{le numéro (30) est déjà celui de la formule de $\mu(G)$, p.~13 de l'auteur ; le « $Y_\alpha/$ » est ajouté au-dessus de la ligne.}
de sorte que $\hat{M}_G$ représente le groupoïde fibré

\page{29}\note{p.~14 de l'auteur.}
\[
\text{(31)}\qquad \mathcal{S} \longmapsto
\begin{array}{l}
\text{groupoïde des systèmes } (Y_\alpha)_{\alpha\in S} \text{ de courbes de MD } Y_\alpha \text{ de genre } g_\alpha\text{,} \\
\text{sur } \mathcal{S}\text{, munies de mono.\ } \hat{I}_{\alpha\,\mathcal{S}} \hookrightarrow Y_\alpha^{\text{lisse}} .
\end{array}
\]
\note{« systèmes $(Y_\alpha)_{\alpha\in S}$ de » et « $Y_\alpha$ de genre $g_\alpha$, » sont ajoutés au-dessus de la ligne ; un mot biffé illisible précède « sur ».}
Posant alors
\[
\text{(32)}\qquad
\begin{cases}
\hat{I} = \coprod \hat{I}_\alpha = I \sqcup \tilde{A} \\
Y = \coprod_{\alpha\in S} Y_\alpha
\end{cases}
\]
\note{devant l'accolade, deux mots biffés, \uncertain{on aura}.}
on aura un mono
\[
\text{(33)}\qquad \hat{I}_{\mathcal{S}} \hookrightarrow Y^{\text{lisse}}
\]
qui permet, d'une part de construire une
\[
\text{(34)}\qquad \mathfrak{X} = Y/(\sigma_{\tilde{A}_{\mathcal{S}}}) ,
\]
qui sera manifestement une courbe \struck{standard} \add{MD} sur $\mathcal{S}$, et d'autre part de définir
\[
\text{(35)}\qquad I_{\mathcal{S}} \hookrightarrow \mathfrak{X}^{\text{lisse}}
\]
\note{dans (34) et (35), et dans la ligne qui suit, des chapeaux sur $\mathfrak{X}$ et sur $I$ sont griffonnés : il semble les supprimer.}
de sorte que $(\mathfrak{X}, I_{\mathcal{S}})$ devient une courbe \struck{standard sur $\mathcal{S}$} MD sur $\mathcal{S}$, « épinglée » par l'ens fini $I$. On a donc un morphisme canonique de multiplicités
\[
\text{(36)}\qquad \hat{M}_G \longrightarrow \hat{M}_{g,I}
\]
\marginal{\emph{NB} Ce morphisme est une imm.\ locale, mais pas une immersion [il \ill{} un morphisme \ill{} $\Gamma = \operatorname{Aut} G$]}
compatible d'ailleurs avec l'action de $\Gamma$ sur $\hat{M}_G$, et son action sur $\hat{M}_{g,I}$ via $\Gamma \to \mathfrak{S}_I$.

\page{30}
On en déduit donc un morphisme canon.\ de multiplicités schématiques

\begin{tikzcd}[column sep=small, row sep=small]
\text{(37)}\quad \hat{M}_{[G]} \arrow[r, hook] \arrow[d, no head, "\wr\ \text{déf}"'] & \hat{M}_{g,\nu} \arrow[d, no head, "\wr"] \\
(\hat{M}_G, \Gamma) & (\hat{M}_{g,I}, \mathfrak{S}_I)
\end{tikzcd}

qui est encore une \emph{immersion}, fermée cette fois-ci, de multiplicités schématiques propres et lisses sur $\mathbb{Z}$. Ainsi on trouve explicitement l'adhérence de $M_{[G]}$ dans $\hat{M}_{g,\nu}$, \struck{comme} sous-multiplicité propre et lisse sur $\operatorname{Spec} \mathbb{Z}$ --- qui \struck{pourrait} \add{\uncertain{espérer}} mieux ! En même temps, cette construction permet de décrire de façon pratiquement triviale les types $[G']$ qui sont spécialisation de $[G]$, puisque ce sont ceux qui proviennent de $\hat{M}_G$. Nous allons \struck{\ill{}} revenir là-dessus maintenant.

\page{31}\note{p.~15 de l'auteur.}
\section{5. Spécialisation \struck{et générisation} de MD-graphes.}

Soit
\[
\text{(38)}\qquad G = (S, \tilde{A}, \sigma_{\tilde{A}}, I, \hat{I} \to S \xrightarrow{g} \mathbb{N}), \qquad \hat{I} = \tilde{A} \sqcup I
\]
\note{la formule est encadrée ; « $\hat{I} = \tilde{A} \sqcup I$ » est écrit sous $\hat{I}$, et un mot surchargé au-dessus de $\sigma_{\tilde{A}}$ est illisible.}
un MD-graphe, de type numérique $(g,\nu)$. Nous allons construire ses \struck{\ill{}} spécialisations, \struck{comme} les \add{MD-}graphes associés aux diverses « strates » de \struck{\ill{}}
\[
\hat{M}_G = \prod_{\alpha\in S} \hat{M}_{g_\alpha, \hat{I}_\alpha} .
\]
Une telle strate est obtenue en choisissant, pour tout $\alpha$, une strate de $\hat{M}_{g_\alpha, \hat{I}_\alpha}$, c'est à dire une classe d'isomorphie de MD-graphes \struck{$G_\alpha$ de type numérique} de type \struck{$(g_\alpha, \hat{\nu}_\alpha)$} numérique $(g_\alpha, \hat{I}_\alpha)$
\[
\text{(39)}\qquad G(\alpha) = \bigl(S(\alpha), \tilde{A}(\alpha), \sigma_{\tilde{A}(\alpha)}, I(\alpha) = \hat{I}_\alpha, \ \hat{I}(\alpha) \overset{\text{déf}}{=} \tilde{A}(\alpha) \sqcup I(\alpha) \to S(\alpha) \xrightarrow{g_\alpha} \mathbb{N}\bigr)
\]
\note{au-dessus de (39), une première écriture de la même formule est biffée, $(S(\alpha), \tilde{A}(\alpha), \sigma_{\tilde{A}(\alpha)}, I(\alpha) = \hat{I}_\alpha \to S(\alpha) \to \mathbb{N})$ ; dans (39) même, la fin de la première écriture est biffée et récrite en dessous, et un groupe biffé précède « : ».}
(\emph{NB} bien entendu, le cas où $G(\alpha)$ est le MD-graphe trivial de type $(g_\alpha, \hat{\nu}_\alpha)$ n'est pas exclu). Par « classes d'isomorphie », il faut entendre ici : pour des isomorphismes induisant l'identité sur $I(\alpha) = \hat{I}_\alpha$. Au lieu de travailler avec de telles classes d'isomorphie, nous allons pour commencer travailler

\page{32}
avec ces $G(\alpha)$ eux-mêmes.

Le \add{MD-}graphe \add{$G'$} spécialisé de $G$ associé aux $G(\alpha)$ s'obtient ainsi :
\[
\text{(40)}\qquad
\left\{
\begin{array}{l}
S' = \coprod_{\alpha\in S} S(\alpha) \\[2pt]
\tilde{A}' = \tilde{A} \sqcup \coprod_{\alpha\in S} \tilde{A}(\alpha) \\[2pt]
\sigma_{\tilde{A}'} \text{ induit par } \sigma_{\tilde{A}} \text{ et par les } \sigma_{\tilde{A}(\alpha)} \\[2pt]
I' = I \quad \text{donc} \quad \hat{I}' \overset{\text{déf}}{=} \tilde{A}' \sqcup I' \simeq \underbrace{(\tilde{A} \sqcup I)}_{\hat{I} = \coprod_{\alpha\in S} \hat{I}_\alpha = \coprod_{\alpha\in S} I(\alpha)} \sqcup \coprod_{\alpha\in S} \tilde{A}(\alpha) \\[2pt]
\phantom{I' = I \quad \text{donc} \quad \hat{I}'} \simeq \coprod_{\alpha\in S} \underbrace{\hat{I}_\alpha \sqcup \tilde{A}(\alpha)}_{\hat{I}(\alpha)} \\[2pt]
\hat{I}' \to S' \text{ induit par les } \hat{I}(\alpha) \to S(\alpha) \\[2pt]
S' \xrightarrow{g'} \mathbb{N} \text{ induit par les } S(\alpha) \xrightarrow{g_\alpha} \mathbb{N}
\end{array}
\right.
\]
\note{deux groupes biffés, devant « $I' = I$ » et au début de la ligne suivante, sont illisibles.}
Donc « la » recette \add{« pratique »} est la suivante :

\struck{\ill{}} On remplace chaque sommet \add{$\alpha \in S$} de $G$ par un MD-graphe $G(\alpha)$, de type numérique $(g_\alpha, \hat{\nu}_\alpha)$, où $\hat{\nu}_\alpha = \nu_\alpha + \tilde{\mu}_\alpha$ ($\nu_\alpha = \operatorname{card} I_\alpha$, poids marqué de $\alpha$ ; $\tilde{\mu}_\alpha$ = ordre de $\alpha$ = nb d'arêtes de $G$ issues de $\alpha$), plus précisément, un MD-graphe de type $(g_\alpha, \hat{I}_\alpha)$, où $\hat{I}_\alpha = I_\alpha \sqcup \tilde{A}_\alpha$.

\page{33}\note{p.~16 de l'auteur.}
Cela signifie qu'on se donne, pour $\alpha$ fixé,

1°) un graphe \add{connexe} $(S(\alpha), \tilde{A}(\alpha), \struck{\ill{}}\ \sigma_{\tilde{A}(\alpha)}, \tilde{A}(\alpha) \to S(\alpha))$

2°) on \struck{\ill{}} se donne une application « genre »
\[
S(\alpha) \xrightarrow{g_\alpha} \mathbb{N}
\]

3°) on donne une application $I_\alpha \to S(\alpha)$, i.e.\ on « répartit » les \struck{\ill{}} « pts marqués » au dessus de $\alpha$ \add{(i.e.\ les éléments de $I_\alpha$)} entre les sommets $\beta \in S(\alpha)$ de $G(\alpha)$

4°) on donne une application $\tilde{A}_\alpha \to S(\alpha)$, i.e.\ pour tt \struck{flèche} \add{arc} du graphe $G$, d'origine $\alpha$, on \struck{se donne} « choisit une \add{nouvelle} origine » parmi les sommets $\beta \in S(\alpha)$ du graphe $G(\alpha)$

5°) les données 1°) : 4°) définissent un \add{nouveau $G(\alpha)$} quadruplet, et on prend soin que celui-ci soit un MD-graphe, \struck{\ill{}} ce qui signifie ici que pour tt $\beta \in S(\alpha)$ « de genre 0 » i.e.\ tel que $g_\alpha(\beta) = 0$, on ait
\[
\operatorname{card} I_{\alpha,\beta} + \operatorname{card} \tilde{A}_{\alpha,\beta} \geqslant 3 .
\]

Ceci étant, on trouve un nouveau \add{MD-}graphe $G'$, \add{qui sera} spécialisé du \add{MD-}graphe $G$.

Je voudrais préciser de façon \uncertain{simplifiée} la relation entre les MD-graphes $G$ et $G'$, et j'introduis à cet effet une digression sur

\page{34}
une \uncertain{catégorisation} de la notion de \emph{morphismes de graphes}.

\struck{Soit}

\section{6. Morphismes de générisation de graphes \add{\ill{}} \struck{ordinaires}}

Soit $G = (S, \tilde{A}, \sigma_{\tilde{A}}, o_{\tilde{A}} : \tilde{A} \to S)$ un graphe ordinaire. On lui associe de façon bien connue un \emph{groupoïde fondamental} \add{$\Pi_1(G)$}, dont l'ens des objets est $S$, et l'ens des \struck{\ill{}} \add{morphismes} de $s$ à $s'$ est l'ens des « chemins simples » de $s$ à $s'$, qui correspondent aux classes d'homotopie de chemins de $s$ à $s'$ dans la réalisation topologique du graphe $G$. Le groupoïde $\Pi_1(G)$ dépend fonctoriellement (dans un sens le plus strict !) du graphe $G$. D'ailleurs, un morphisme de graphes
\[
\text{(41)}\qquad f : G' \to G''
\]

\page{35}\note{p.~17 de l'auteur.}
est connu quand on connait
\[
\text{(42)}\qquad \Pi_1(f) : \Pi_1(G') \to \Pi_1(G'') ,
\]
puisqu'alors on connait l'action de $f$ sur $S$ via l'action de $\Pi_1(f)$ sur les objets, et l'action de $f$ sur les arcs via l'action de $\Pi_1(f)$ sur les flèches, grâce au diagramme commutatif

\begin{tikzcd}
\text{(43)}\quad \tilde{A}(G') \arrow[r, "\tilde{A}(f)"] \arrow[d, hook] & \tilde{A}(G'') \arrow[d, hook] \\
\mathrm{Fl}(\Pi_1(G')) \arrow[r, "{\mathrm{Fl}(\Pi_1(f))}"] & \mathrm{Fl}(\Pi_1(G''))
\end{tikzcd}

En fait, les morphismes de $G'$ dans $G''$ s'identifient aux morphismes de groupoïdes $\Pi_1(G')$ dans $\Pi_1(G'')$, qui transforment arcs \add{de $G'$} en arcs de $G''$.

Ceci étant, on généralise la notion de morphisme de graphes en celle de \emph{morphisme généralisé} de $G'$ dans $G''$ (ou « \emph{morphisme de générisation} » de $G'$ dans $G''$)

\page{36}
comme étant les morphismes de groupoïdes $\Pi_1(G') \to \Pi_1(G'')$ qui transforment tout arc de $G'$ en une flèche de $\Pi_1(G'')$ qui est, \emph{soit} un arc de $G''$, soit une flèche identique. En d'autres termes, notant \emph{flèches élémentaires} de $\Pi_1(G')$ les flèches qui sont, soit identiques, soit un \struck{arc} arc de $G'$ (ensemble \add{de flèches} stable par passage à la flèche opposée…), les morphismes \struck{de} générisation de $G'$ dans $G''$ sont les morphismes de $\Pi_1(G')$ dans $\Pi_1(G'')$ qui transforment flèches élémentaires en fl.\ élémentaires. Cette description a l'avantage de rendre évident que cette notion est stable par composition.

Pratiquement, la donnée d'un morphisme \struck{de} générisation \struck{de} $G' \to G''$ est définie

\page{37}\note{p.~18 de l'auteur.}
par une application
\[
\text{(44)}\qquad S' \sqcup \tilde{A}' \longrightarrow S'' \sqcup \tilde{A}''
\]
qui applique $S'$ dans $S''$ (mais pas néc.\ $\tilde{A}'$ dans $\tilde{A}''$), qui commute à $\sigma_{G'}$ et $\sigma_{G''}$ (définis via $\sigma_{\tilde{A}'}$ et $\sigma_{\tilde{A}''}$, avec l'identité sur $S'$ resp.\ $S''$) et commute à l'application « origine » (qui sur $S'$ resp.\ $S''$ est l'identité). On peut aussi dire que c'est défini par la donnée

\marginal{(45) et il faut supposer de plus que pour tt $a \in \tilde{A}' \setminus \tilde{B}'$, l'origine et l'extrémité de $a$ ont \ill{} dans $S$ par $f_S$}
(45) 1°) d'une application $f_{S} : S' \to S''$

2°) d'une partie \struck{$\tilde{A}''$} de $\tilde{A}'$, stable par $\sigma_{\tilde{A}'}$\note{la partie est notée $\tilde{B}'$ à l'alinéa suivant.}

3°) d'une application $\tilde{B}' \to \tilde{A}''$ commutant à $\sigma_{\tilde{A}'} | \tilde{B}'$ et à $\sigma_{\tilde{A}''}$, et aux applications $o_{\tilde{A}'} | \tilde{B}'$ et $o_{\tilde{A}''}$.

On peut dire aussi que c'est aussi la donnée d'un sous-graphe $G'_{\tilde{B}'}$ de \struck{\ill{}} $G'$ ayant mêmes sommets (cela défini justement par la donnée d'une partie $\tilde{B}'$ de $\tilde{A}'$ stable par $\sigma_{\tilde{A}'}$), et un morphisme de $G'_{\tilde{B}'}$ dans $G''$ ! La composition des morphismes généralisés est donc définie de façon

\page{38}
évidente. Ainsi on trouve la catégorie \struck{\ill{}} des graphes, dont les objets sont les graphes, les \struck{morphismes} \add{flèches} les morphismes généralisés de graphes. On notera que les iso.\ de cette nouvelle catégorie sont les iso au sens ordinaire.

Supposons maintenant que \add{$G'$, $G''$} $\underline{G}'$, $\underline{G}''$ fassent partie de MD-graphes, i.e.\ soient munis de plus de \struck{\ill{}}
\[
\text{(46)}\qquad
\left|
\begin{array}{l}
S \xrightarrow{g} \mathbb{N}, \quad S' \xrightarrow{g'} \mathbb{N} \\
I \to S, \quad I' \to S'
\end{array}
\right.
\]
\note{il écrit $G'$, $G''$ au-dessus de $\underline{G}'$, $\underline{G}''$ ; (46) porte $S$, $S'$ et non $S'$, $S''$.}
On dit qu'on a un \emph{morphisme de générisation} de $\underline{G}'$ dans $\underline{G}''$, si on se donne, en plus des morphismes \add{généralisés} $G' \to G''$, une application bijective
\[
\text{(47)}\qquad I' \simeq I'' ,
\]
\struck{satisfai} et si les conditions suivantes sont satisfaites :

\page{39}\note{p.~19 de l'auteur. Le but du morphisme s'appelle désormais $G$, $S$, $\tilde{A}$, $I$ (et non plus $G''$, …) ; la page est très surchargée.}
a) $S' \to S''$ surjectif, $\tilde{B}' \to \tilde{A}''$ bijectif

b) $I' \simeq I$, $I' \to S'$, $I \to S$, $S' \to S$ commutatif :

\begin{tikzcd}
I' \arrow[r, "\sim"] \arrow[d] & I \arrow[d] \\
S' \arrow[r] & S
\end{tikzcd}

\marginal{donc $\tilde{B}' \sqcup I' \xrightarrow{\sim} \tilde{A} \sqcup I \to S$ et $\tilde{B}' \sqcup I' \to S' \to S$ commutent.}

c) $\forall \alpha \in S$, le sous-graphe $G'(\alpha)$ de $G'$ [formé des \struck{\ill{}} $s' \in S'$ au-dessus de $\alpha$, et des \struck{flèches} arcs $a' \in \tilde{A}'$ qui sont dans \add{$\tilde{C}' =$} $\tilde{A}' \setminus \tilde{B}'$ (i.e.\ au-dessus de $\mathrm{id}_\alpha$) et dont l'origine (donc l'extrémité aussi) est dans $S'_\alpha$] muni de \struck{\ill{}} \add{$\tilde{B}' \sqcup$} $I'_\alpha \hookrightarrow S'_\alpha$ induit par \add{$\tilde{B}' \sqcup$} $I' \to S'$, et de $g_\alpha : S'_\alpha \to \mathbb{N}$ induit par $g' : S' \to \mathbb{N}$, est un \emph{MD-graphe} \add{$\underline{G}'(\alpha)$}, i.e.\ il est \emph{connexe}, et [il satisfait \ill{}] à la condition de rigidité relative aux sommets $\beta$ tels que $g'_\beta \in \{0,1\}$, à savoir que
\[
2g'_\beta + (\operatorname{card} I'_\beta + \operatorname{card} \tilde{B}'_\beta) \geqslant 3
\]
\note{les additions et suppressions autour de cette condition sont nombreuses et en partie illisibles : un « $+ \operatorname{card} \tilde{B}'_\beta$ » est ajouté au-dessus de la ligne, une flèche renvoie $\tilde{B}'_\beta$ à « ordre de $\beta$ dans $\underline{G}'(\alpha)$ », et à droite il récrit « i.e.\ : $2g'_\beta + \operatorname{card} I'_\beta + \operatorname{card} \tilde{B}'_\beta$ » ; un bloc de deux lignes en dessous, biffé, portait des variantes de la même condition, $2g'_\beta + \operatorname{card} I'_\beta \geqslant 3$ et $2g'_\beta + \operatorname{card} I'_\beta + \operatorname{card} \tilde{A}'_\beta \geqslant 3$, avec « une condition plus faible … puisque $\underline{G}'$ est un MD-graphe ». Après « MD-graphe », un passage biffé de deux lignes (« [NB il \ill{}] ») est illisible.}
\marginal{\emph{NB} $\tilde{B}'_\alpha$ est l'ens des $a' \in \tilde{B}'$ dont l'origine se trouve dans $S'_\alpha$, \ill{} donc $\tilde{B}'_\alpha \xrightarrow{\sim} \tilde{A}_\alpha$}
\marginal{poids marqué en $\beta$ \ill{} dans le \ill{} graphe $\underline{G}'(\alpha)$}

d) Le genre du MD-graphe $\underline{G}(\alpha)$ est égal à $g_\alpha$, i.e.\ on a
\[
\sum_{\beta\in S'_\alpha} g'_\beta + h_1(G(\alpha)) = g_\alpha .
\]
\note{il écrit $G(\alpha)$ dans d), là où c) définit $G'(\alpha)$ ; cf.\ la condition 3°) de la page suivante, qui porte $h_1(G'(\alpha))$.}

\page{40}
\emph{NB} Pour qu'un morphisme généralisé de graphes $G' \to G$ corresponde à un morphisme de générisation (pour $I' \simeq I$ convenable), il f\supplied{au}t et s\supplied{uffi}t qu'il satisfasse aux conditions suivantes

(48)

1°) $S' \to S$ surjectif, $\tilde{B}' \to \tilde{A}$ bijectif, les « graphes fibres » $G'(\alpha)$ ($\alpha \in S$) connexes. (Ces conditions ne dépendent pas des MD-structures sur $G'$, $G$)

2°) $\forall \alpha \in S$, on a
\[
\underbrace{\operatorname{card} I_\alpha}_{\nu_\alpha} = \sum_{\beta\in S'_\alpha} \underbrace{\operatorname{card} I'_\beta}_{\nu'_\beta}
\]
(cette condition ne dépend que des « poids marqués » sur $S$, $S'$, et bien sûr de $S' \to S$)

3°) $\forall \alpha \in S$, on a
\[
g_\alpha = \Bigl(\sum_{\beta\in S'_\alpha} g'_\beta\Bigr) + h_1(G'(\alpha))
\]

\struck{4°) $\forall \alpha' \in S'$, on a (non seulement}
\struck{$2g'_{\alpha'} + \operatorname{card} I'_{\alpha'}$ $+ \operatorname{card} \tilde{A}'_{\beta'}$ $\geqslant 3$, mais aussi)}
\struck{$2g'_\alpha + \operatorname{card} I'_{\alpha'}$ $+ \operatorname{card} \tilde{B}'_{\alpha'}$ $\geqslant 3$}
\note{la condition 4°) est barrée de grandes croix. Les conditions (48) sont réunies par une accolade. L'énumération se poursuit au-delà du lot.}


\end{document}
